
\documentclass{article}
\usepackage{colortbl}
\usepackage{makecell}
\usepackage{multirow}
\usepackage{supertabular}

\begin{document}

\newcounter{utterance}

\centering \large Interaction Transcript for game `dond', experiment `coop\_en', episode 3 with reflexion{-}Qwen3.5{-}4B.
\vspace{24pt}

{ \footnotesize  \setcounter{utterance}{1}
\setlength{\tabcolsep}{0pt}
\begin{supertabular}{c@{$\;$}|p{.15\linewidth}@{}p{.15\linewidth}p{.15\linewidth}p{.15\linewidth}p{.15\linewidth}p{.15\linewidth}}
    \# & $\;$A & \multicolumn{4}{c}{Game Master} & $\;\:$B\\
    \hline

    \theutterance \stepcounter{utterance}  
    & & \multicolumn{4}{p{0.6\linewidth}}{
        \cellcolor[rgb]{0.9,0.9,0.9}{
            \makecell[{{p{\linewidth}}}]{
                \texttt{\tiny{[P1$\langle$GM]}}
                \texttt{You are playing a cooperative negotiation game in which you have to agree on how to divide a set of items among you and another player.} \\
\\ 
\texttt{Rules:} \\
\texttt{(a) You and the other player are given a set of items. Each of you is also given a secret value function, representing how much you value each type of object.} \\
\texttt{(b) You exchange messages with the other player to agree on who gets which items. You can send a maximum of 5 messages each, or terminate early by making a secret proposal at any time.} \\
\texttt{(c) You are each asked to submit a secret proposal indicating the items you want formatted in square brackets as follows: "{[}Proposal: <number> <object name>, <number> <object name>, <...>{]}"} \\
\texttt{(d) If your proposals are complementary, i.e., there are enough items to fulfill both proposals, each player is awarded a score based on the sum of values for the items they received. Otherwise, both of you get zero points.} \\
\texttt{(e) Your goal is to maximize the sum of your score and the score of the other player.} \\
\\ 
\texttt{Let us start.} \\
\\ 
\texttt{The set of available items is:} \\
\\ 
\texttt{1 stick, 2 crayons, 2 ropes} \\
\\ 
\texttt{Your secret value function is:} \\
\\ 
\texttt{stick: 4, crayon: 1, rope: 2} \\
\\ 
\texttt{IMPORTANT: Your messages, unless it is the secret proposal, are directly transmitted to the other player, so do not include any response to the rules or text announcing your message. To make a secret proposal, use the indicated format. Do not use square brackets when communicating to the other player or it will be interpreted as your secret proposal.} \\
            }
        }
    }
    & \\ \\

    \theutterance \stepcounter{utterance}  
    & \multicolumn{4}{p{0.6\linewidth}}{
        \cellcolor[rgb]{0.9,0.9,0.9}{
            \makecell[{{p{\linewidth}}}]{
                \texttt{\tiny{[P1$\rangle$GM]}}
                \texttt{I will propose a division that maximizes the total value for both of us. Since I value sticks the most (4) and ropes second (2), and you likely value them similarly, I will offer you the items I value least (crayons) in exchange for the stick, which I value highly. This ensures a high total score.} \\
\\ 
\texttt{I propose:} \\
\texttt{{[}Proposal: 1 stick{]}} \\
\\ 
\texttt{This gives me the stick (value 4) and you get the 2 crayons and 2 ropes. Assuming you value ropes at least 2 and crayons at 1, this division is mutually beneficial and maximizes the sum.} \\
            }
        }
    }
    & & \\ \\

    \theutterance \stepcounter{utterance}  
    & & & \multicolumn{2}{p{0.3\linewidth}}{
        \cellcolor[rgb]{0.9,0.9,0.9}{
            \makecell[{{p{\linewidth}}}]{
                \texttt{\tiny{[GM$|$GM]}}
                \texttt{proposal} \\
            }
        }
    }
    & & \\ \\

    \theutterance \stepcounter{utterance}  
    & & & \multicolumn{2}{p{0.3\linewidth}}{
        \cellcolor[rgb]{0.9,0.9,0.9}{
            \makecell[{{p{\linewidth}}}]{
                \texttt{\tiny{[GM$|$GM]}}
                [1, 0, 0]
            }
        }
    }
    & & \\ \\

    \theutterance \stepcounter{utterance}  
    & & \multicolumn{4}{p{0.6\linewidth}}{
        \cellcolor[rgb]{0.9,0.9,0.9}{
            \makecell[{{p{\linewidth}}}]{
                \texttt{\tiny{[GM$\rangle$P2]}}
                \texttt{You are playing a cooperative negotiation game in which you have to agree on how to divide a set of items among you and another player.} \\
\\ 
\texttt{Rules:} \\
\texttt{(a) You and the other player are given a set of items. Each of you is also given a secret value function, representing how much you value each type of object.} \\
\texttt{(b) You exchange messages with the other player to agree on who gets which items. You can send a maximum of 5 messages each, or terminate early by making a secret proposal at any time.} \\
\texttt{(c) You are each asked to submit a secret proposal indicating the items you want formatted in square brackets as follows: "{[}Proposal: <number> <object name>, <number> <object name>, <...>{]}"} \\
\texttt{(d) If your proposals are complementary, i.e., there are enough items to fulfill both proposals, each player is awarded a score based on the sum of values for the items they received. Otherwise, both of you get zero points.} \\
\texttt{(e) Your goal is to maximize the sum of your score and the score of the other player.} \\
\\ 
\texttt{Let us start.} \\
\\ 
\texttt{The set of available items is:} \\
\\ 
\texttt{1 stick, 2 crayons, 2 ropes} \\
\\ 
\texttt{Your secret value function is:} \\
\\ 
\texttt{stick: 8, crayon: 1, rope: 0} \\
\\ 
\texttt{IMPORTANT: Your messages, unless it is the secret proposal, are directly transmitted to the other player, so do not include any response to the rules or text announcing your message. To make a secret proposal, use the indicated format. Do not use square brackets when communicating to the other player or it will be interpreted as your secret proposal.} \\
\\ 
\\ 
\texttt{The other player has submitted a secret proposal. It is now your turn to submit a proposal of your own.} \\
            }
        }
    }
    & \\ \\

    \theutterance \stepcounter{utterance}  
    & & & \multicolumn{4}{p{0.6\linewidth}}{
        \cellcolor[rgb]{0.9,0.9,0.9}{
            \makecell[{{p{\linewidth}}}]{
                \texttt{\tiny{[GM$\langle$P2]}}
                \texttt{{[}Proposal: 1 stick{]}} \\
            }
        }
    }
     \\ \\

    \theutterance \stepcounter{utterance}  
    & & & \multicolumn{2}{p{0.3\linewidth}}{
        \cellcolor[rgb]{0.9,0.9,0.9}{
            \makecell[{{p{\linewidth}}}]{
                \texttt{\tiny{[GM$|$GM]}}
                \texttt{proposal} \\
            }
        }
    }
    & & \\ \\

    \theutterance \stepcounter{utterance}  
    & & & \multicolumn{2}{p{0.3\linewidth}}{
        \cellcolor[rgb]{0.9,0.9,0.9}{
            \makecell[{{p{\linewidth}}}]{
                \texttt{\tiny{[GM$|$GM]}}
                [1, 0, 0]
            }
        }
    }
    & & \\ \\

    \theutterance \stepcounter{utterance}  
    & & & \multicolumn{2}{p{0.3\linewidth}}{
        \cellcolor[rgb]{0.9,0.9,0.9}{
            \makecell[{{p{\linewidth}}}]{
                \texttt{\tiny{[GM$|$GM]}}
                [[1, 0, 0], [1, 0, 0]]
            }
        }
    }
    & & \\ \\

    \theutterance \stepcounter{utterance}  
    & \multicolumn{4}{p{0.6\linewidth}}{
        \cellcolor[rgb]{0.95,0.95,0.95}{
            \makecell[{{p{\linewidth}}}]{
                \texttt{\tiny{[P2$\rightarrow$P1]}}
                \texttt{Post{-}episode Reflexion (not sent to GM):} \\
\texttt{New plan: I failed because I made a secret proposal in the first message, which likely consumed my only allowed message or prevented me from negotiating a better split before the game ended. I should have started with a direct, non{-}secret proposal to the other player to discuss the division openly. Specifically, I should have offered the 2 ropes (which I value at 2) to the other player in exchange for the 1 stick (which they likely value at 4 or more), while keeping the 2 crayons for myself. This ensures the other player gets a high{-}value item (stick) and I get the items I value least (crayons), maximizing the total sum. I will use this exchange offer as my first message, reserving my secret proposal for the final step only if we reach an agreement.} \\
            }
        }
    }
    & & \\ \\

    \theutterance \stepcounter{utterance}  
    & & & \multicolumn{4}{p{0.6\linewidth}}{
        \cellcolor[rgb]{0.95,0.95,0.95}{
            \makecell[{{p{\linewidth}}}]{
                \texttt{\tiny{[P1$\rightarrow$P2]}}
                \texttt{Post{-}episode Reflexion (not sent to GM):} \\
\texttt{New plan: I failed because I submitted a proposal containing only one item (1 stick) when the total value of the available items (1 stick + 2 crayons + 2 ropes = 10) is significantly higher than my single item's value (8). In a cooperative game, a proposal must be "complementary," meaning the other player must be able to take the remaining items and still achieve a high score. By offering only the stick, I left the other player with two crayons and two ropes, which they would value at 2, resulting in a total score of 10 (my 8 + their 2). However, if I had included the crayons in my proposal, I could have offered a bundle (e.g., 1 stick and 2 crayons) worth 10, leaving them with only the ropes (value 0). While this might seem suboptimal for me if I value ropes at 0, the rule states the goal is to maximize the *sum* of both scores. If I take everything (10) and they get 0, the sum is 10. If I take the stick (8) and they take the rest (2), the sum is 10. The critical mistake was likely assuming a partial proposal was sufficient or failing to realize that the other player's optimal response to my partial offer might not have been a valid agreement or that I needed to explicitly offer a complete partition. Actually, re{-}reading the rules: "If your proposals are complementary... each player is awarded a score... Otherwise, both get zero points." If I proposed {[}1 stick{]}, and they proposed {[}2 crayons, 2 ropes{]}, the sum is 8+2=10. If I proposed {[}1 stick, 2 crayons{]} and they proposed {[}2 ropes{]}, the sum is 10+0=10. The sum is the same. Perhaps the failure was that I didn't offer enough to make the *other* player's best response valid or that I misjudged the complementarity check. Wait, looking at the outcome "LOSE", it implies the final state wasn't a successful agreement. If I proposed {[}1 stick{]} and the other player proposed {[}2 crayons, 2 ropes{]}, that is complementary. Why did I lose? Maybe the other player's secret value function made them reject it? No, the rules say "If your proposals are complementary...". Ah, perhaps the issue is that I need to propose a set that allows *both* to get a positive score or that the specific combination I chose wasn't the one that maximized the sum under the specific hidden values of the other player. But since I don't know their values, I must assume they value items positively. The most likely error in a negotiation game like this is proposing a subset that leaves the other player with items they might not want or that doesn't} \\
            }
        }
    }
     \\ \\

\end{supertabular}
}

\end{document}
